% Generated draft: requires mathematical and source review.
\documentclass[11pt]{article}
\usepackage[margin=25mm]{geometry}
\usepackage{fontspec}
\setmainfont{Latin Modern Roman}
\usepackage{amsmath,amssymb,mathtools,unicode-math}
\setmathfont{Latin Modern Math}
\usepackage{enumitem,needspace,xurl,hyperref}
\setlength{\parindent}{0pt}
\setlength{\parskip}{5pt}
\setlength{\emergencystretch}{3em}
\newcommand{\N}{\mathbb{N}}
\newcommand{\Z}{\mathbb{Z}}
\newcommand{\Q}{\mathbb{Q}}
\newcommand{\R}{\mathbb{R}}
\newcommand{\C}{\mathbb{C}}
\newcommand{\F}{\mathbb{F}}
\newcommand{\A}{\mathbb{A}}
\newcommand{\PP}{\mathbb{P}}
\newcommand{\E}{\mathbb{E}}
\newcommand{\eps}{\varepsilon}
\newcommand{\dd}{\,\mathrm{d}}
\begin{document}
\section*{UnsolvedMath ID: 700867}
\textbf{Source problem code:} catalogue entry
\subsection{Definitions and mathematical statement}
Let \(\Omega\) be a set with at least two elements, and let \(G\leq\operatorname{Sym}(\Omega)\). The action is \emph{sharply doubly transitive} if, for all \(\alpha,\beta,\gamma,\delta\in\Omega\) with \(\alpha\neq\beta\) and \(\gamma\neq\delta\), there exists exactly one \(g\in G\) such that \(g\alpha=\gamma\) and \(g\beta=\delta\). Fix \(\omega\in\Omega\) and write
\[
H=G_\omega=\{g\in G:g\omega=\omega\}.
\]
A group is \emph{locally finite} if each of its finitely generated subgroups is finite. A subgroup \(N\leq G\) acts \emph{regularly} on \(\Omega\) if, for every \(\alpha,\beta\in\Omega\), there exists exactly one \(n\in N\) such that \(n\alpha=\beta\).

For every sharply doubly transitive permutation group as above, ask the following two questions separately:
\begin{enumerate}
\item[(a)] If \(H\) is locally finite, must there exist a normal subgroup \(N\trianglelefteq G\) acting regularly on \(\Omega\)?
\item[(b)] If there exists an abelian subgroup \(A\leq H\) with \([H:A]<\infty\), must there exist a normal subgroup \(N\trianglelefteq G\) acting regularly on \(\Omega\)?
\end{enumerate}
No finiteness assumption on \(G\) or \(\Omega\) is imposed. Since point stabilizers are conjugate, either hypothesis may equivalently be imposed on any point stabilizer.

\emph{Supplied qualification (Comments of 2022).} The source reports an affirmative answer to part~(b) in permutation characteristic \(0\), attributed to F. O. Wagner's 2022 preprint. In the standard convention for sharply doubly transitive groups, permutation characteristic \(0\) means that involutions (elements of order two) have a fixed point and the product of any two distinct involutions has infinite order. The reported affirmative answer is restricted to this characteristic and is supplied source information, not an independently verified result here.
\subsection{Short English statement}
Must a sharply doubly transitive permutation group have a regular normal subgroup whenever a point stabilizer is locally finite, or whenever a point stabilizer contains an abelian subgroup of finite index? These are separate questions. The supplied 2022 comment reports an affirmative answer to the second question in permutation characteristic \(0\).
\subsection{Sources}
\begin{itemize}
\item Local UnsolvedMath registry, permanent ID \(700867\), exact source problem code \texttt{catalogue entry}, title \emph{Kourovka Notebook Problem 12.48}. The supplied registry labels the problem open; this label is not an independent verification of its mathematical status.
\item V. D. Mazurov, proposer, 1992. \emph{The Kourovka Notebook (Unsolved Problems in Group Theory)}, Problem 12.48, originating in the 12th issue (1992), as recorded in the supplied 21st-issue edition (2026), \texttt{arXiv:1401.0300v48}. The supplied collection description names E. I. Khukhro and V. D. Mazurov as editors.
\item The source refers to Archive 11.52 for the definition of sharply doubly transitive; that archive entry is not included in the supplied record.
\item Comments of 2022: F. O. Wagner, \emph{Preprint}, 2022, \url{https://hal.archives-ouvertes.fr/hal-03590818}. The supplied comment attributes an affirmative answer to part~(b) in permutation characteristic \(0\) to this preprint; its contents have not been independently checked here.
\end{itemize}
\end{document}
